\section{Introduction}

Post-trained language models produce far less varied output than their base counterparts, and the
loss is large enough to undermine inference-time scaling, synthetic data generation, and any
application that depends on sampling a population of responses rather than a single one
\citep{karouzos2026where,matthews2026language,yun2025price}. The phenomenon is now documented
across formats \citep{parikh2026structured}, personas \citep{xiao2026chameleons}, and successive
post-training stages \citep{liu2026representation}, and remedies have been proposed at the level
of training objectives \citep{li2025jointly} and sampling procedures \citep{xu2026prism}.

A second line of work treats the problem as one of representation rather than optimisation.
Because activation-space interventions can reliably shift sentiment, sycophancy, refusal, and
truthfulness \citep{turner2023steering,rimsky2024steering,zhao2025steering}, it is natural to
ask whether diversity has its own direction: contrast the generations of one prompt that came out
varied with those that came out repetitive, average the activation difference, and add the
resulting vector during inference. The recipe inherits the standard contrastive-activation
assumption---that a behaviour of interest is encoded along a single linear direction recoverable
from labelled pairs---and it is attractive precisely because it needs no retraining and no human
annotation. Recent work has begun to question how far that assumption travels, finding that
steering vectors generalise poorly across contexts \citep{tan2024analyzing}, that residual-stream
addition can fail outright in stateful dialogue \citep{kang2026promptactivation}, that the
contrastive pairs determine the outcome more than the fitting procedure
\citep{bohek2026detecting,hao2025patterns}, and that linear, context-independent interventions are
a simplification that new methods must move beyond \citep{jafari2026metasteer,park2026information}.
None of this work examines diversity, where the label is not given but must be constructed by the
experimenter from the model's own samples.

That construction is the weak point. When the positive and negative classes are not annotated
categories but quantiles of a dispersion score computed over the same generations being compared,
the contrast direction is as much a statistic of sampling noise as of the model. We fit the
direction exactly as the recipe prescribes and then split the evidence in two. When the halves come
from different prompts the fit is almost perfectly reproducible, cosine $+0.999$ for pythia-410m;
when they come from the \emph{same} prompt it is unreproducible, $-0.037$ for the same model. The
pooled estimator is stable because prompts have stable identities, not because diversity has a
stable direction, so averaging across prompts yields a sum of near-orthogonal, individually
unpredictable components: a quantity with a norm but no direction.

We turn that diagnosis into a set of predictions and test each one. If the pooled vector is not a
diversity axis, then (i) injecting it should not change diversity at any dose; (ii) it should not
outperform a random vector of the same norm; (iii) fitting a direction per prompt, which removes
the identity contamination, should still fail because the within-prompt signal is itself unstable;
and (iv) removing the response subspace should nevertheless collapse diversity, because ablation
destroys generative capacity whether or not the removed subspace is semantically special. All four
predictions hold. Interventions at 35 dose levels across seven models yield no positive
dose--response in any model, and the fitted direction is numerically worse than a random one at
matched coefficients. Prompt-local steering succeeds no better than steering with a different
prompt's direction. Ablation does collapse diversity---to 1.35 \vendi{} from a baseline of 6.13 in
pythia-1b---but a control that keeps the activations and destroys only the high/low pairing
collapses it just as far, to 1.46, while a random direction leaves it at 6.05.

Our contributions are as follows.
\begin{itemize}
\item The standard contrastive construction of a ``diversity direction'' fails a split-half test at
the within-prompt level in all seven models we study while passing it at the pooled level, locating
its apparent stability in prompt identity rather than in diversity (\Cref{sec:stability}).
\item A full placebo ladder---random, pairing-shuffled, orthogonalised, and centroid
directions---with per-model dose calibration against a labelled positive control finds no
diversity-increasing effect anywhere and a significantly negative median slope
(\Cref{sec:interventions}).
\item The only significant placebo effects vanish once generations that shorten or repeat are
excluded, so they measure degeneration rather than control (\Cref{sec:interventions}).
\item Diversity is removable along a subspace but not addable along a direction, and the removal
does not depend on the contrast that defines the direction (\Cref{sec:mechanism}).
\end{itemize}
Together these results argue that diversity should not be treated as another linearly steerable
attribute, and they specify the controls that any future claim of diversity steering should report.
